\documentclass[11pt]{article}
\usepackage[margin=1in]{geometry}
\usepackage{amsmath,amssymb,amsthm,mathtools,microtype}
\usepackage[hidelinks]{hyperref}
\newtheorem{theorem}{Theorem}[section]
\newtheorem{proposition}[theorem]{Proposition}
\newtheorem{lemma}[theorem]{Lemma}
\newtheorem{corollary}[theorem]{Corollary}
\theoremstyle{remark}\newtheorem{remark}[theorem]{Remark}
\newcommand{\E}{\mathbb E}
\newcommand{\R}{\mathbb R}
\newcommand{\G}{\mathcal G}
\newcommand{\F}{\mathcal F}
\newcommand{\norm}[1]{\left\lVert#1\right\rVert}
\title{Gaussian likelihood transport:\\square-integrable domains, past projections\\and exact innovation energy}
\author{Project Stardust\\\small A classical formalization companion}
\date{October 29, 2026}
\begin{document}\maketitle
\begin{abstract}
Equivalent probability measures preserve null sets but need not preserve square-integrable domains. We reconstruct an explicit delayed Gaussian likelihood for which the terminal density is normalized despite failure of the Novikov exponential moment, and its square root belongs to the original $L^2$ space but not the tilted one. Multiplication by the square root is nevertheless a unitary between the correct spaces. The entire time-one past has the same law under both measures, while its transported Hilbert subspace differs from the original past subspace. We prove the correct conditional projection formula, an exact innovation energy $1-\sqrt{2/3}$ and agreement after the likelihood becomes observable. Further elementary deductions give the exact operator-norm distance one, an instability under arbitrarily small likelihood-root changes, and an explicit two-stage predictable coefficient. These deductions are distinguished from named formal declarations. The mechanisms are classical; no new change-of-measure or martingale representation theorem is claimed.
\end{abstract}

\section{Spaces and the precise comparison}
Let $(\Omega,\F,P)$ carry a standard Brownian motion $W$, with raw natural past
$\F_t=\sigma(W_s:0\le s\le t)$. All spaces here are real Hilbert spaces of equivalence classes. The verified realization uses the original countable Gaussian product and its constructed Brownian-generated subspace; it does not identify the entire ambient Gaussian space with the Brownian-generated space. The calculations below use actual finite Gaussian laws and independence of an increment from the whole past.

Put
\begin{equation}\label{eq:XY}
 X=W_1,\qquad Y=W_2-W_1,\qquad
 Z=\exp(-XY-X^2/2),\qquad R=\sqrt Z.
\end{equation}
Under $P$, $(X,Y)$ has the product of two standard Gaussian laws. Write
$\gamma(dx)=\phi(x)\,dx$, where
$\phi(x)=(2\pi)^{-1/2}e^{-x^2/2}$. We will prove $\E_P Z=1$, so $dQ=Z\,dP$ is a probability law equivalent to $P$.

For a sub-$\sigma$-algebra $\G$, let $H_P(\G)=L^2(\Omega,\G,P)$ and similarly define $H_Q(\G)$. These are closed subspaces of the corresponding ambient $L^2$ spaces, with measurable versions understood up to the ambient null sets. Define
\[
 U:L^2(Q)\longrightarrow L^2(P),\quad Uf=Rf,
 \qquad V(\G)=U H_Q(\G).
\]
Let $P_\G$ be the ordinary $P$-conditional projection and $\Pi_\G$ the orthogonal projection onto $V(\G)$. The central comparison is between two operators on the \emph{same} space $L^2(P)$, namely $P_\G$ and $\Pi_\G$. Comparing conditional expectations before specifying this identification would conceal the domain issue.

The source map in the accompanying artifact fixes every named formal correspondence to the frozen verified baseline. General subspace criteria, parameter extensions, the exact operator-norm result and the explicit running coefficient below are complete human deductions from these classical structures, not additional Lean theorems.

\section{A normalized likelihood beyond the Novikov moment}
For real $a,y$, set
\[
 k_a(y)=e^{-ay-a^2/2},\qquad r_a(y)=e^{-ay/2-a^2/4}.
\]
Completing a square gives the pointwise identity
\begin{equation}\label{eq:shift}
 \phi(y)k_a(y)=\phi(y+a).
\end{equation}
In particular $k_a>0$, $\int k_a\,d\gamma=1$, and $r_a^2=k_a$.

\begin{proposition}\label{prop:normal}
The variable $Z$ in \eqref{eq:XY} is positive, integrable and normalized. However
\[
 \E_P Z^2=\infty,\qquad \E_P e^{X^2/2}=\infty.
\]
\end{proposition}
\begin{proof}
Tonelli and \eqref{eq:shift} give
$\E Z=\int\int k_x(y)\,d\gamma(y)d\gamma(x)=1$. Positivity then gives integrability. Also
\[
 k_x(y)^2=e^{x^2}k_{2x}(y),\qquad
 \E Z^2=\int e^{x^2}\,d\gamma(x)=\infty.
\]
For the last assertion, $e^{x^2/2}\phi(x)=(2\pi)^{-1/2}$ on the whole real line, whose Lebesgue integral is infinite. The preceding integral with $e^{x^2}$ is larger. All infinite-moment statements concern nonnegative integrals, so no undefined signed cancellation occurs.
\end{proof}

There is an actual predictable integrand behind this density. Let
\[
 \theta_t=X\,\mathbf1_{(1,2]}(t).
\]
Its coefficient is measurable at the left endpoint, and
\[
 \int_0^2\theta_t^2\,dt=X^2,\qquad
 \E\int_0^2\theta_t^2\,dt=1,\qquad
 \int_0^2\theta_t\,dW_t=XY.
\]
Thus $Z=\exp(-\int\theta\,dW-\frac12\int\theta^2dt)$ is a normalized terminal stochastic exponential with infinite Novikov expectation. This is failure of a sufficient test, not failure of the terminal likelihood.

For completeness, the associated density process can be checked without invoking that test. It is one up to time one, equals
\[
 Z_t=\exp\{-X(W_t-W_1)-(t-1)X^2/2\}\quad(1\le t\le2),
\]
and is constant afterwards. Conditionally on $\F_1$, each future increment is Gaussian; its exponential moment cancels the displayed compensator. The same conditioning at any two times in $[1,2]$ proves the martingale identity. Conditioning across time one uses $\E[Z_t\mid\F_1]=1$ and the tower property. Consequently $Z_t=\E[Z\mid\F_t]$ throughout this finite horizon. This direct process calculation is a classical deduction in the exposition; the formal source explicitly constructs the step integral, terminal density, normalization and failed exponential moment.

\section{Actual past law and covariance under the tilt}
Let $\mathsf H=(W_s)_{0\le s\le1}$ as a random element of the coordinate product space of histories, with law $\mu$. Independence of future Brownian increments says that $(\mathsf H,Y)$ has law $\mu\otimes\gamma$. The endpoint evaluation $x(h)=h(1)$ is measurable.

\begin{theorem}\label{thm:past}
Under $Q$, the pair $(\mathsf H,W_2)$ has law $\mu\otimes\gamma$. Thus $Q$ and $P$ agree on every event of $\F_1$, and $W_2$ is independent of the entire time-one past under $Q$.
\end{theorem}
\begin{proof}
For any nonnegative measurable $f$ of a history and a real coordinate, change variables $z=y+x(h)$ in the inner integral to obtain
\begin{align*}
 \E_Q f(\mathsf H,W_2)
 &=\int\!\int k_{x(h)}(y)f(h,y+x(h))\,d\gamma(y)d\mu(h)\\
 &=\int\!\int f(h,z)\,d\gamma(z)d\mu(h).
\end{align*}
The equality holds first for nonnegative functions by Tonelli and then for integrable functions by their positive and negative parts. Taking $f$ to be a history indicator proves equality of all past event probabilities. Product indicators prove the claimed independence, with no restriction to finitely many observed times.
\end{proof}

\begin{corollary}\label{cor:schur}
Under $Q$, the covariance and inverse precision of $(X,Y)$ are respectively
\[
 C=\begin{pmatrix}1&-1\\-1&2\end{pmatrix},\qquad
 A=\begin{pmatrix}2&1\\1&1\end{pmatrix},\qquad AC=CA=I.
\]
Their quadratic energy has the two exact forms
\begin{equation}\label{eq:schur}
 2x^2+2xy+y^2=x^2+(y+x)^2
 =2(x+y/2)^2+y^2/2.
\end{equation}
\end{corollary}
\begin{proof}
By the theorem, $X$ and $S=W_2=X+Y$ are independent standard normals under $Q$. Hence $Y=S-X$ has variance two and covariance $-1$ with $X$. This gives $C$; ordinary matrix multiplication gives its stated inverse, which also follows from the density
$\phi(x)\phi(y)k_x(y)=\phi(x)\phi(y+x)$. Expanding squares gives \eqref{eq:schur}. Minimizing over $y$ leaves $x^2$, while minimizing over $x$ leaves $y^2/2$. The corresponding Schur values are one and one half; neither value is an unchanged copy of both diagonal entries.
\end{proof}

In particular, the original coordinate process $W$ is not a $Q$-martingale even at this one step: $\E_Q[W_2\mid\F_1]=0$, whereas $W_1=X$ is nonzero almost surely. This does not concern a corrected process with its drift removed.

\section{The exact square-integrable domains}
\begin{theorem}\label{thm:domain}
The measures $P,Q$ have the same null sets, but
\[
 R\in L^2(P),\quad \norm{R}_{L^2(P)}=1,
 \qquad R\notin L^2(Q).
\]
No change of representative in the original $P$-class repairs the latter failure. In contrast, $U f=Rf$ is a surjective linear isometry from $L^2(Q)$ to $L^2(P)$, with inverse $v\mapsto v/R$.
\end{theorem}
\begin{proof}
Strict positivity and finiteness of $Z$ imply both absolute continuities: if $P(B)=0$, then $Q(B)=\int_B Z\,dP=0$; conversely, $Q(B)=0$ forces $Z\mathbf1_B=0$ almost surely and hence $P(B)=0$. Now
\[
 \E_P R^2=\E_P Z=1,\qquad
 \E_Q R^2=\E_P Z^2=\infty.
\]
Equality of null ideals transfers equality of representatives from $P$ to $Q$, so it cannot change the second integral.

For every measurable $f$, including when the value is infinite,
\[
 \int |Rf|^2\,dP=\int |f|^2\,dQ.
\]
This proves the precise domain equivalence $Rf\in L^2(P)\iff f\in L^2(Q)$, independence of representatives, linearity and preservation of norms. For $v\in L^2(P)$, positivity gives
$\int |v/R|^2dQ=\int |v|^2dP$, and multiplication and division cancel almost surely. Thus the inverse is defined on the whole target space. Polarization proves inner-product preservation.
\end{proof}

The distinction is between an unchanged formula and a transported class: $R$ fails to belong to $L^2(Q)$, while $U^{-1}R=1$ belongs to $L^2(Q)$. Equivalence of measures alone supplies no bounded identity map between their full $L^2$ spaces. On the other hand, Theorem~\ref{thm:past} gives exactly the same $L^2$ norm under both laws for every $\F_1$-measurable representative, because their restricted measures agree. The domain obstruction therefore concerns the terminal class and is compatible with equality of all past norms.

\section{Conditional transport and a criterion for agreement}
The following general statement applies to any positive normalized density $Z$; Gaussian structure is unnecessary. Write $R=\sqrt Z$, $Z_\G=\E_P[Z\mid\G]$ and define $U,V(\G)$ as in Section 1.

\begin{theorem}\label{thm:projection}
The image $V(\G)$ is closed, and its orthogonal projection satisfies
\begin{equation}\label{eq:projection}
 \Pi_\G v=R\E_Q[v/R\mid\G]
 =\frac{R}{Z_\G}\E_P[Rv\mid\G],\qquad v\in L^2(P).
\end{equation}
Moreover
\begin{equation}\label{eq:criterion}
 V(\G)=H_P(\G)\quad\Longleftrightarrow\quad
 R\text{ has a }\G\text{-measurable version}.
\end{equation}
\end{theorem}
\begin{proof}
The measurable subspace $H_Q(\G)$ is closed: from an $L^2$-convergent sequence one can choose an almost surely convergent subsequence, whose limit has a measurable version. A unitary maps a closed subspace to a closed subspace. Conjugating its orthogonal projection therefore gives the first formula.

The variable $v/R$ is $Q$-integrable by Cauchy--Schwarz, and $Rv$ is $P$-integrable by the same inequality. Also $Z_\G>0$ almost surely: its zero set $B\in\G$ has $\int_B Z\,dP=0$, and positivity of $Z$ implies $P(B)=0$. The conditional Bayes formula gives the second expression. To verify that formula without any boundedness of $1/Z_\G$, apply the defining integral identity on sets in $\G$ with $Z_\G$ bounded away from zero and infinity, then exhaust their union. Nonnegative truncations justify pull-out and integrability; in particular
\[
 \E_Q\left|\frac{\E_P[Rv\mid\G]}{Z_\G}\right|
 =\E_P\left|\E_P[Rv\mid\G]\right|\le\E_P|Rv|.
\]
Testing against $\G$-indicators establishes the desired conditional expectation.

If the subspaces agree, $1\in H_Q(\G)$ implies $R=U1\in H_P(\G)$. Conversely, if $R$ has a $\G$-measurable positive version, multiplication by it takes each $\G$-measurable $Q$-class into a $\G$-measurable $P$-class. Division by it gives the reverse inclusion with exactly the integrability established in Theorem~\ref{thm:domain}. This proves \eqref{eq:criterion}.
\end{proof}

Equality in \eqref{eq:criterion} is also equivalent to equality of the two orthogonal projections, since a projection determines its range. The general equivalence and Bayes rewriting are explicit classical deductions. The formal source supplies the actual unitary, image subspace and first conditional formula in this example. It also proves equality with the ordinary past after time two, when $R$ is already measurable.

\section{Full-past projection error and conditional variance}
We now return to \eqref{eq:XY}, and write $P_1=P_{\F_1}$ and $\Pi_1=\Pi_{\F_1}$. Conditionally on the whole past, $X$ is fixed and $Y$ is standard Gaussian. The exponential moment of $Y$ gives
\begin{equation}\label{eq:means}
 \E_P[R\mid\F_1]=e^{-X^2/8}=:m(X),\qquad
 \E_P[R^2\mid\F_1]=1.
\end{equation}
Indeed $\int r_x(y)d\gamma(y)=e^{-x^2/4}e^{x^2/8}$, and the second identity is normalization of $k_x$. These formulas can equivalently be verified by integration over arbitrary measurable history sets, so conditioning on $X$ has not replaced conditioning on the whole past.

\begin{theorem}\label{thm:energy}
The actual conditional variance is almost surely
\[
 \operatorname{Var}_P(R\mid\F_1)=1-e^{-X^2/4}>0.
\]
Furthermore
\begin{equation}\label{eq:error}
 \norm{P_1R}_2^2=\sqrt{2/3},\qquad
 \norm{R-P_1R}_2^2=1-\sqrt{2/3},\qquad \Pi_1 R=R.
\end{equation}
Consequently the ordinary and transported past subspaces and projections differ.
\end{theorem}
\begin{proof}
Subtracting the square of the conditional mean in \eqref{eq:means} gives the variance. A standard normal is nonzero almost surely, so the exponential is strictly below one almost surely. The Gaussian square integral
\begin{equation}\label{eq:square}
 \int e^{-c x^2}\,d\gamma(x)=(1+2c)^{-1/2},\qquad c\ge0,
\end{equation}
follows by the substitution $u=\sqrt{1+2c}\,x$ in the normalized Gaussian integral. At $c=1/4$ it gives the first norm in \eqref{eq:error}. Orthogonality of conditional expectation gives the second. Finally $R=U1$ belongs to every transported past, so its projection is itself. The positive error separates the two projections and the two closed ranges.
\end{proof}

This yields the formally recorded lower bound
\[
 1-\sqrt{2/3}\le\norm{\Pi_1-P_1}_{\mathrm{op}}^2,
\]
since $R$ is a unit vector. There is no contradiction with $Q|_{\F_1}=P|_{\F_1}$: the vectors in the image $U H_Q(\F_1)$ include a terminal multiplier that need not be past-measurable under $P$.

\section{Exact norm distance and a parameter instability}
The previous lower bound is not the exact operator norm. The following sharper statements are human deductions from the same conditional fibers.

\begin{theorem}\label{thm:normone}
For the above example,
\[
 \norm{\Pi_1-P_1}_{\mathrm{op}}=1.
\]
More generally, for any real $a$, put
\[
 Z_a=e^{-aXY-a^2X^2/2},\quad R_a=\sqrt{Z_a},\quad dQ_a=Z_a\,dP,
\]
and let $\Pi_{1,a}$ be the projection transported from $L^2(Q_a,\F_1)$. Then
\begin{equation}\label{eq:instability}
 \norm{\Pi_{1,a}-P_1}_{\mathrm{op}}=1\quad(a\ne0),\qquad
 \Pi_{1,0}=P_1,
\end{equation}
although $\norm{R_a-1}_2\to0$ as $a\to0$.
\end{theorem}
\begin{proof}
Normalization, equivalence, preservation of the whole past law and unitary transport follow as before with $x$ replaced by $ax$. In particular
\[
 \E[R_a^2\mid\F_1]=1,\qquad
 \E[R_a\mid\F_1]=m_a(X):=e^{-a^2X^2/8}.
\]
Fix $a\ne0$ and $M>0$. The event $B_M=\{|X|>M\}$ has positive probability. The vector
\[
 v_M=R_a\mathbf1_{B_M}/\sqrt{P(B_M)}
\]
has norm one and lies in the transported past. Its ordinary past projection is $m_a(X)\mathbf1_{B_M}/\sqrt{P(B_M)}$. Hence
\[
 \norm{(\Pi_{1,a}-P_1)v_M}_2^2
 =\frac{\E[(1-m_a(X)^2)\mathbf1_{B_M}]}{P(B_M)}
 \ge1-e^{-a^2M^2/4}.
\]
Taking $M\to\infty$ gives the lower bound one. For any two orthogonal projections $E,F$, the vectors $E(I-F)v$ and $(I-E)Fv$ are orthogonal, and their difference is $(E-F)v$. Therefore
\[
 \norm{(E-F)v}^2
 =\norm{E(I-F)v}^2+\norm{(I-E)Fv}^2
 \le\norm{(I-F)v}^2+\norm{Fv}^2=\norm{v}^2.
\]
This proves the matching upper bound.

For the final assertion, all $R_a$ have norm one, and \eqref{eq:square} gives
\[
 \E R_a=(1+a^2/4)^{-1/2},\qquad
 \norm{R_a-1}_2^2=2\{1-(1+a^2/4)^{-1/2}\}\longrightarrow0.
\]
At $a=0$, both measures and the unitary are the identity case, proving \eqref{eq:instability}.
\end{proof}

The witnesses concentrate on rare past events depending on $a$ and $M$. Thus a fixed-vector likelihood-root error does not bound a worst-case projection error. No conclusion about failure of strong-operator convergence follows from this norm calculation. It compares the transported projection with $P_1$ on the common $L^2(P)$ space; it does not compare two untransported operators with different domains.

\section{Likelihood-root overlap and conditional geometry}
The preceding calculation also has an exact two-parameter form. For real $u,v$, direct multiplication gives
\[
 r_u(y)r_v(y)=e^{-(u-v)^2/8}k_{(u+v)/2}(y).
\]
All integrals are finite, since each square has integral one. Integrating yields the classical Gaussian square-root overlap
\begin{equation}\label{eq:overlap}
 \int r_u r_v\,d\gamma=e^{-(u-v)^2/8},\qquad
 \int(r_u-r_v)^2d\gamma=2\{1-e^{-(u-v)^2/8}\}.
\end{equation}
In particular the second quantity is positive when $u\ne v$. The source proves these exact fiber formulas.

For the actual parameter family above, use $u=aX,v=bX$ in \eqref{eq:overlap}, condition on the whole past, and apply \eqref{eq:square}. This proves the further human deduction
\[
 \langle R_a,R_b\rangle_{L^2(P)}
 =\{1+(a-b)^2/4\}^{-1/2},\qquad
 \norm{R_a-R_b}_2^2=2\left\{1-\frac{1}{\sqrt{1+(a-b)^2/4}}\right\}.
\]
These are square-root-density distances in one common reference space. They give continuous dependence of each likelihood-root vector on its parameter, while Theorem~\ref{thm:normone} exhibits the separate operator-norm phenomenon. No statistical estimation theorem or infinite-dimensional Gaussian equivalence criterion is inferred from this finite calculation.

\section{An explicit predictable coefficient and its energy}
The formal source proves that $R$ is a genuine Brownian-generated terminal $L^2(P)$ value, with a unique original predictable representer supported in $[0,2]$. Its future energy after time one is exactly the error in \eqref{eq:error}. Here we give a direct classical construction of that representer, avoiding any additional derivative-domain assumption on $R$.

\begin{proposition}\label{prop:integrand}
Let $M_t=\E_P[R\mid\F_t]$ for $0\le t\le2$. Versions are
\begin{equation}\label{eq:Mt}
 M_t=\begin{cases}
 \displaystyle\frac{2}{\sqrt{5-t}}\exp\!\left(-\frac{W_t^2}{2(5-t)}\right),&0\le t\le1,\\[6pt]
 \displaystyle\exp\!\left(-\frac{X(W_t-W_1)}2-\frac{tX^2}{8}\right),&1\le t\le2.
 \end{cases}
\end{equation}
Define $h_t$ on these two open time intervals by
\begin{equation}\label{eq:ht}
 h_t=\begin{cases}
 -W_t M_t/(5-t),&0<t\le1,\\
 -X M_t/2,&1<t\le2.
 \end{cases}
\end{equation}
Endpoint choices at single deterministic times are immaterial. Then $h$ is predictable, belongs to $L^2(dt\otimes P)$, and
\begin{align}
 R-2/\sqrt5&=\int_0^2h_t\,dW_t,\label{eq:representation}\\
 \E\int_0^2h_t^2dt&=1/5,\qquad
 \E\int_1^2h_t^2dt=1-\sqrt{2/3}.\label{eq:energies}
\end{align}
\end{proposition}
\begin{proof}
For $1\le t\le2$, write $Y=(W_t-W_1)+(W_2-W_t)$. The last increment is independent of $\F_t$, with variance $2-t$. Its Gaussian exponential moment gives the second expression in \eqref{eq:Mt}. At time one this is $e^{-X^2/8}$. For $t\le1$, condition this time-one expression on $W_t$, using the independent remaining increment of variance $1-t$. Completing the square in its Gaussian integral gives
\[
 \E\left[e^{-(w+\sqrt{1-t}\,G)^2/8}\right]
 =\{1+(1-t)/4\}^{-1/2}
   e^{-w^2/[8\{1+(1-t)/4\}]},
\]
which is the first expression. The formulas agree at time one; at time two the value is $R$, and at zero it is $2/\sqrt5$. They define a continuous adapted version, so the coefficients in \eqref{eq:ht} are predictable on their intervals.

For $f(t,w)=2(5-t)^{-1/2}e^{-w^2/[2(5-t)]}$, ordinary differentiation gives
\[
 f_w=-\frac{w}{5-t}f,\qquad f_t+\tfrac12f_{ww}=0.
\]
For fixed $X=x$, the second function in \eqref{eq:Mt}, viewed as a function of $t$ and $W_t-W_1$, has spatial derivative $-xf/2$ and again solves the backward heat equation. It\^o's formula on each interval, localized to bounded paths and bounded coefficients, therefore gives $dM_t=h_t\,dW_t$ with no time drift. The random parameter $X$ is known at time one and remains constant on the second interval; it is not differentiated as a running state variable there.

Here is the integrability step for removing localization. Let $\tau_n$ be increasing bounded stopping times tending to two, localizing $M$ and $h$. Since $M_t=\E[R\mid\F_t]$, Jensen and optional conditioning give $\E M_{\tau_n}^2\le\E R^2=1$. The stopped It\^o isometry consequently bounds
\[
 \E\int_0^{\tau_n}h_t^2dt
 =\E(M_{\tau_n}-M_0)^2\le1-M_0^2=1/5.
\]
One may justify optional conditioning here by approximating each bounded stopping time from above by finite-valued times: continuity of $M$ and the uniform integrability of its squares, inherited from conditional expectations of $R^2$, give the limit. Monotone convergence now puts $h$ in the full square-integrable predictable domain. Its stopped integrals converge by the isometry, and continuity at two identifies the limit with $R-M_0$. This proves \eqref{eq:representation} and the first energy identity. The same isometry from one to two gives
$\E\int_1^2h_t^2dt=\E(R-M_1)^2$, evaluated in Theorem~\ref{thm:energy}.
\end{proof}

The energy up to one is therefore $\sqrt{2/3}-4/5>0$. Uniqueness in the original predictable $L^2$ domain follows from the isometry: two coefficients with the same terminal integral have difference of zero $L^2(dt\otimes P)$ norm. This explicitly constructed coefficient agrees with the formally supplied unique predictable representer. We do not identify it with a projected Malliavin derivative without separately proving membership in that derivative's closed domain. Representation by the bounded adjoint of the It\^o isometry does not require such a membership claim.

\section{Observable-time agreement and scope}
For $t\ge2$, $R$ is $\F_t$-measurable. Theorem~\ref{thm:projection} therefore gives
\[
 U H_Q(\F_t)=H_P(\F_t),\qquad
 P_t U=U\,\E_Q[\,\cdot\mid\F_t],\qquad t\ge2.
\]
The unitary is always valid, but its commutation with an ordinary past projection depends on when the terminal multiplier becomes observable. The failure at time one and the agreement after time two are both statements about actual closed Hilbert subspaces, not an assertion that raw filtrations under equivalent laws have different event sets.

The frozen sources formally supply the normalized delayed step, actual full-past law, covariance/precision calculation, $L^2$ counterexample, half-density unitary, image projections, exact one-vector error/conditional variance, future predictable energy and eventual agreement. The general criterion, exact operator norm, parameter family and explicit running coefficient are proved classical deductions in this exposition. Every formal source is unchanged. No $Q$-Brownian construction, transported Malliavin graph, general Girsanov theorem, observed-IBM representation or new stochastic calculus is claimed.

\section{Classical context and attribution}
Conditional Bayes transport and Gaussian completion of squares are classical. Allaart's author-hosted notes give the ordinary Gaussian change of measure and the general conditional Bayes formula, Lemma 5.10 \cite{allaart}; the argument above includes its needed integrability details. Ruf's primary paper states the classical Novikov sufficiency condition and develops much more general martingality criteria \cite{ruf}. Our normalized delayed example is not a new substitute for those criteria. Gaussian square-root overlap is likewise an elementary classical distance calculation.

The bounded comparison did not establish priority for the exact assembled numerical example or for the operator-norm instability formulation. This companion's purpose is complete domain-sensitive exposition and traceability to actual verified source declarations. None of the classical constructions, or the elementary deductions printed here, is advertised as a new general probability theorem. The exact source map and separate finite algebra checks document the reconstruction; they do not replace the analytic proofs or independent referee review.

\begin{thebibliography}{9}
\bibitem{allaart} P. Allaart, \emph{Math 6810 (Probability), Chapter 5 lecture notes}, University of North Texas, Fall 2012, dated February 19, 2013, Section 5.3 and Lemma 5.10. \url{https://itservices.cas.unt.edu/~allaart/6810/6810-notes-chap5.pdf}.
\bibitem{ruf} J. Ruf, \emph{A new proof for the conditions of Novikov and Kazamaki}, Stochastic Processes and their Applications \textbf{123} (2013), 404--421; primary preprint \href{https://arxiv.org/abs/1111.5583}{arXiv:1111.5583}, introduction and Corollary 1.
\end{thebibliography}
\end{document}
